\Needspace{14\baselineskip}
\section{A configuration with two coordinatizations}
\label{sec:gadget}

We now use the deletion lemma to construct the exceptional configuration.
We need two properties of the same set: many points on each undeleted
line, and rigidity of its small line covers.  A change of coordinates
lets us obtain both.  In the original coordinates we delete one line
in every finite nonzero direction.  After the change of coordinates,
the deleted sets become curves, so the line-survival conclusion of
Lemma~\ref{lem:linecover} applies to every line.  It will force every
small mixed cover by lines and individual points to be a parallel class.

\begin{lemma}
\label{lem:gadget}
For every sufficiently large prime \(s\equiv2\pmod3\), there is
a set \(Y\subseteq\F_s^2\) with the following properties.  For each
\(d\in\F_s^*\), there is an affine line \(\ell_d\) of slope \(d\),
and
\[
 Y=\F_s^2\setminus\bigcup_{d\in\F_s^*}\ell_d.
\]
For each \(d\in\F_s^*\), every line of slope \(d\) other than
\(\ell_d\) contains at least five points of \(Y\).  The map
\[
 f:\F_s^2\longrightarrow\F_s^2,
 \qquad f(x,y)=(x^3,y),
\]
is a bijection, and every affine line contains at least \(s/8\) points
of \(f(Y)\).  Moreover, if \(a\) affine lines and \(b\) individual
points cover \(f(Y)\), and
\begin{equation}
\label{con-weighted-budget}
 a+\frac b2\le s,
\end{equation}
then \(a=s\), \(b=0\), and the lines form a full parallel class.
Consequently every such mixed cover has weighted cost at least \(s\),
and no cover of \(f(Y)\) by at most \(s\) lines contains two
nonparallel lines.
\end{lemma}

\begin{proof}
Let \(s\) tend to infinity through primes congruent to \(2\pmod3\).
Since
\(\gcd(3,s-1)=1\), cubing permutes \(\F_s^*\) and fixes zero.
Thus \(f\) is bijective.  It maps horizontal lines to horizontal
lines and vertical lines to vertical lines.  We first prove the two
survival properties simultaneously with high probability, then use
them to classify mixed covers satisfying~\eqref{con-weighted-budget}.

\paragraph{Survivors in the original coordinates.}
Choose the lines \(\ell_d\), independently for the \(s-1\) values
of \(d\), uniformly within their parallel classes.  Fix a line
\(\ell\) of finite nonzero slope \(d\), and temporarily ignore the
choice \(\ell_d\).  Each of the other \(s-2\) choices deletes a
uniformly distributed point of \(\ell\).  Let \(Z_\ell\) count
the surviving points on \(\ell\).  Its expectation is
\[
 \mu:=\E Z_\ell=s(1-1/s)^{s-2}.
\]
Changing one of these choices changes the survivor count by at most
two.  Apply the independent-trial jump bound and bounded-drift
concentration estimate of
\cite[Lemmas~A.2--A.3]{OpenAI2026Balanced} to the exposure martingale,
with \(T=s-2\), \(L=2\), and \(u=\mu/2\).  The martingale has
zero drift, and hence
\[
 \Pr(Z_\ell<\mu/2)
 \le 2\exp\left(-\frac{\mu^2}{128(s-2)}\right)
 =\exp(-\Omega(s)).
\]
There are fewer than \(s^2\) lines under consideration.  A union
bound, together with \(\mu/2\ge5\) for sufficiently large \(s\),
shows that all these counts are at least five with probability tending
to one.  If \(\ell\ne\ell_d\), reinstating the deletion by
\(\ell_d\) changes no point on \(\ell\), since the two lines are
parallel.  Thus every required line in the original coordinates has
at least five points of \(Y\).

\paragraph{Survivors and almost-pencil covers in the new coordinates.}
In the new coordinates, the possible deletion shapes are
\[
 F_{d,z}=\{(x^3,dx+z):x\in\F_s\},
 \qquad d\in\F_s^*,\quad z\in\F_s.
\]
There are \(s(s-1)\) such shapes, each of size \(s\), and each of
the \(s-1\) independent trials chooses one of them.  Every point
is deleted by a specified trial with probability \(1/s\), so its
overall survival probability is the uniform value
\[
 \delta_s=(1-1/s)^{s-1}\ge\frac14
\]
for all sufficiently large \(s\).

Each shape meets every affine line in at most three points.  For a
vertical line this follows from bijectivity of cubing; for a
horizontal line it follows from \(d\ne0\).  For a line of finite
nonzero slope \(m\), the intersection is determined by
\[
 dx+z=mx^3+u,
\]
a nonzero polynomial equation of degree three.  In particular, no
shape is itself a line when \(s>5\).  All the hypotheses of
Lemma~\ref{lem:linecover} hold with field size \(s\), a fixed
integer constant, for example \(D=3\), survival lower bound
\(\delta=1/4\), and \(\rho_0=1/100\).  The exact-line
occurrence probabilities are zero, and the subfield restriction is
automatic for a prime field.  With probability tending to one, every
line consequently contains at least \(s/8\) points of \(f(Y)\),
and its almost-cover conclusion holds for \(f(Y)\) as well.  We now
combine that conclusion with the lower bound on every line to replace
an almost pencil by a complete one.

\paragraph{From an almost pencil to a full parallel class.}
Assume these events hold, and consider a cover satisfying
\eqref{con-weighted-budget}.  Repeated resources can be removed;
alternatively, one may keep their original counts, since removal
only reduces the budget.  The line portion has at most \(s\)
distinct lines and leaves at most \(b\le2s\) points of \(f(Y)\)
uncovered.  Lemma~\ref{lem:linecover} supplies at least
\((1-\rho_0)s\) distinct lines in this portion through a common
projective point \(p\).  There are at most \(\rho_0s\) other
lines, and
\[
 b\le2(s-a)\le2\rho_0s.
\]
If an affine line through \(p\) were absent from the cover, all the
chosen lines through \(p\) would cover at most one of its affine
points.  The other lines would cover at most \(\rho_0s\) further
points, and the individual point resources at most \(2\rho_0s\).
Thus at most \(1+3\rho_0s\) points of this missing line would be
covered.  For sufficiently large \(s\), this contradicts
\[
 \frac s8>1+3\rho_0s.
\]
Every affine line through \(p\) must therefore be present.  A
finite \(p\) requires \(s+1\) lines, which is impossible.  An
infinite \(p\) requires all \(s\) lines in its parallel class,
and then \eqref{con-weighted-budget} forces \(a=s\) and \(b=0\).
This also rules out repetitions in a cover attaining the budget.

The events used in the two coordinate systems hold simultaneously
with probability tending to one.  Consequently every sufficiently
large prime \(s\equiv2\pmod3\) admits a successful configuration
\(Y\).  A mixed cover of cost less than
\(s\) would satisfy \eqref{con-weighted-budget} but not its
necessary equality conclusion.  The two stated consequences follow.
\end{proof}

There are arbitrarily large primes congruent to \(2\pmod3\).  Indeed,
if the list were finite and its product were \(P\), the integer
\(3P-1\equiv2\pmod3\) would have a prime factor congruent to
\(2\pmod3\), and none of the primes dividing \(P\) could divide it.
Fix any sufficiently large prime \(s\equiv2\pmod3\), and choose
\(Y\) and the omitted lines from Lemma~\ref{lem:gadget}.
These data remain fixed for the rest of the proof.

\section{Ambient placement}
\label{sec:placement}

We will place copies of \(Y\) in a larger affine plane, assign each
copy the directions in which its labels will change, and delete the
ambient lines of those directions that meet it.  The placement must
meet two requirements.  Joins between different copies must avoid
their changed directions, while the undeleted points must retain the
pencil and line-survival properties of Lemma~\ref{lem:linecover}.
We first specify the copies and direction sets, then prove the
deletion estimates for independent pools of translates, and finally
select one compatible translate from each pool.

\subsection{Fields, copies, and direction sets}

Let \(n\ge3\) be odd, and put \(q=s^n\). We identify \(\F_s\)
with the prime subfield of \(\F_q\). If \(E\) is a subfield, then
\(|E|=s^a\), where \(a=[E:\F_s]\), and multiplication of vector-space
bases gives
\[
 n=a[\F_q:E].
\]
For a proper subfield, the index \([\F_q:E]\) is an odd integer
larger than one, so it is at least three. Consequently
\[
 |E|=s^a\le s^{n/3}=q^{1/3}
 \qquad\text{for every proper subfield }E.
\]
Thus Theorem~\ref{thm:incidence} and
Lemma~\ref{lem:linecover} apply along this sequence of ambient
fields.  In this section all constants may depend on the already
fixed \(s\) and \(Y\), but not on \(n\).

The odd-degree condition has now supplied the required subfield bound.
It remains to assign every ambient direction to at least one copy,
with bounded overlap between assignments.  Identify the ambient
directions with
\(\PP^1(\F_q)=\F_q\cup\{\infty\}\), using slope.  Choose a
set \(\Lambda\) of representatives of
\(\F_q^*/\F_s^*\), with \(1\in\Lambda\).  There will be one
configuration group for each \(\lambda\in\Lambda\), together
with one further group denoted by \(*\).  Write
\[
 \mathcal I=\Lambda\sqcup\{*\},
 \qquad N_q=|\mathcal I|=\frac{q-1}{s-1}+1.
\]
Each index \(i\in\mathcal I\) will specify a pool of translates
of one configuration.  We call this pool a group and will choose one
translate from each group.

For \(\lambda\in\Lambda\), define
\[
 M_\lambda(x,y)=(x,\lambda y),
 \qquad S_\lambda=\lambda\F_s^*,
 \qquad A_\lambda=\{0,\infty\}.
\]
Choose distinct \(u,v\in\F_s^*\) and define
\[
 M_*=
 \begin{pmatrix}1&1\\u&v\end{pmatrix},
 \qquad
 S_*=(\F_s\cup\{\infty\})\setminus\{u,v\},
 \qquad A_*=\{u,v\}.
\]
All these linear maps are invertible.  For \(M_*\), the image of
a finite slope \(t\) is
\[
 \frac{u+vt}{1+t},
\]
where a zero denominator gives the infinite direction.  Its images
of the horizontal and vertical directions are \(u\) and \(v\).
The induced projective bijection therefore maps the \(s-1\)
finite nonzero subfield slopes precisely to \(S_*\).  The same
statement for \(M_\lambda\) follows from the slope map
\(t\mapsto\lambda t\).

For each \(i\in\mathcal I\), put
\[
 W_i=M_i(\F_s^2).
\]
The directions of joins of distinct points of \(W_i\) form exactly
\(S_i\cup A_i\).  Each set \(S_i\) has size \(s-1\), and
\(A_i\) consists of the two images of the coordinate axes.  The
cosets \(S_\lambda\) partition all nonzero finite ambient slopes.
The extra set \(S_*\) contains \(0\) and \(\infty\).  Hence
every ambient direction belongs to at least one and at most two
sets \(S_i\).  More precisely, the directions belonging to two
sets are
\begin{equation}
\label{con-double-directions}
 S_1\cap S_* = \F_s^*\setminus\{u,v\},
\end{equation}
and the two groups involved are always \(1\) and \(*\).

The sets \(S_i\) are the directions whose labels will change, and
\(A_i\) are the images of the two axis directions preserved by \(f\).
For compatibility between copies we use the
slightly larger direction sets
\[
 D_i=
 \begin{cases}
 S_i\cup A_i,&i\in\{1,*\},\\
 S_i,&i\notin\{1,*\}.
 \end{cases}
\]
In particular,
\begin{equation}
\label{con-conflict-directions}
 D_1=D_*=\PP^1(\F_s),
\end{equation}
whereas the remaining \(D_i\) are the nonidentity multiplicative
cosets of \(\F_s^*\).  Every direction occurs in at most two
sets \(D_i\).  Since \(|W_i|=s^2\), \(|D_i|\le s+1\), and
\(N_q=O_s(q)\), a fixed constant depending only on \(s\)
satisfies all size and multiplicity assumptions of
Theorem~\ref{thm:selection}.  In a conflict-free selection, including
the axes in both \(D_1\) and \(D_*\) ensures that an axis line
meeting one of these two selected copies cannot also meet the other.
This will enter the weighted cover
argument for the two groups sharing a direction
in~\eqref{con-double-directions}.

\subsection{Independent candidate pools}

Let \(K\) be a fixed pool size supplied by
Theorem~\ref{thm:selection}.  For each
\(i\in\mathcal I\), choose independent uniform shifts
\(b_{i,1},\ldots,b_{i,K}\in\F_q^2\), independently between
groups, and write
\[
 P_{i,k}=b_{i,k}+W_i,
 \qquad Y_{i,k}=b_{i,k}+M_iY.
\]
Associate to this candidate the following set of ambient deletion
lines:
\[
 \mathcal B_{i,k}
 =\{\ell:\operatorname{dir}(\ell)\in S_i,
                  \ \ell\cap Y_{i,k}\ne\varnothing\}.
\]

\begin{lemma}
\label{con-deletion-count}
For every candidate and every \(d\in S_i\), exactly \(s-1\)
lines of \(\mathcal B_{i,k}\) have direction \(d\).  Each of
these lines contains at least five points of \(Y_{i,k}\).
Consequently \(|\mathcal B_{i,k}|=(s-1)^2\).
\end{lemma}

\begin{proof}
Undo the translation and apply \(M_i^{-1}\).  An ambient line in
the given direction meeting \(P_{i,k}\) becomes a line through a
point of \(\F_s^2\) in a finite nonzero subfield direction.  Its
intersection with \(\F_s^2\) is the corresponding subfield
line.  The \(s\) lines in that subfield parallel class extend to
distinct parallel ambient lines.  Exactly one is the omitted line
in the definition of \(Y\); the other \(s-1\) each contain at
least five points of \(Y\), by Lemma~\ref{lem:gadget}.  These are
precisely the lines meeting the candidate's retained configuration.
\end{proof}

Each candidate therefore prescribes a bounded list of deletion lines.
We next apply the deletion lemma to all these independent lists at
once.  Besides the almost-pencil and line-survival conclusions, we
need one bound on how many groups a line can meet along joins from an
external point.  All three properties will survive the subsequent
selection of one candidate per group.

\begin{lemma}
\label{con-pool-properties}
There are constants \(\delta\in(0,1)\) and
\(\rho=\delta/10\), independent of \(n\), with the following
properties with probability tending to one as \(n\) tends to
infinity through odd integers. Let
\[
 X_{\rm pool}
 =\F_q^2\setminus
   \bigcup_{i\in\mathcal I}\bigcup_{k=1}^K
          \bigcup_{\ell\in\mathcal B_{i,k}}\ell.
\]
\begin{enumerate}
\item Every cover of \(X_{\rm pool}\) by at most \(q\) affine
lines contains at least \((1-\rho)q\) lines through a common
projective point.
\item For every affine line \(\ell\), if only exact occurrences
of \(\ell\) are ignored in the pool deletions, at least
\(\delta q/2\) points of \(\ell\) remain.
\item For every finite point \(p\) and affine line \(\ell\)
not through \(p\), at most five groups \(i\) have a candidate
\(P_{i,k}\) containing a point \(y\in\ell\) whose join to
\(p\) has direction in \(S_i\).
\end{enumerate}
\end{lemma}

\begin{proof}
We verify the hypotheses of Lemma~\ref{lem:linecover} for the full
pool, which gives the first two assertions.  The third follows from
a separate union bound over finite points and lines.

Treat each candidate shift as one independent deletion trial, listing
the lines of \(\mathcal B_{i,k}\).  There are \(KN_q=O_{s,K}(q)\)
trials, each listing \(m=(s-1)^2\) lines.  We may take the shape
family to be all \(q(q+1)\) ambient affine lines.  Each shape has
size \(q\), and two distinct line shapes intersect in at most one
point.  Thus the shape and trial bounds in
Lemma~\ref{lem:linecover} hold with a sufficiently large fixed
integer \(D\) depending only on \(s,K\).

A candidate's deleting set is a uniform translate of a fixed union
of \(m\) lines.  Its deletion probability at a fixed point is
independent of the point and is at most \(m/q\).  Independence of
the shifts makes the overall survival probability \(\delta_q\)
uniform on \(\F_q^2\).  For sufficiently large \(q\),
\(m/q\le1/2\) and \(N_q\le2q/(s-1)\).  Using
\(\log(1-t)\ge-2t\) for \(0\le t\le1/2\), we obtain
\[
 \delta_q\ge (1-m/q)^{KN_q}
 \ge \exp\bigl(-2KN_qm/q\bigr)
 \ge \exp\bigl(-4K(s-1)\bigr).
\]
Fix
\(\delta=\exp(-4K(s-1))\) and \(\rho=\delta/10\).

For completeness, the exact-line hypothesis also has a uniform
bound.  A fixed line \(\ell\) of direction \(d\) can occur only
in candidates of groups with \(d\in S_i\), of which there are at
most two.  In each such candidate, there are \(s-1\) translated
lines in this direction.  Each fixed line has probability \(1/q\)
of becoming \(\ell\) under the uniform shift.  Therefore the sum,
over all trials, of the probabilities of listing \(\ell\) is at
most
\begin{equation}
\label{con-exact-line-probability}
 \frac{2K(s-1)}q.
\end{equation}
Increasing the fixed \(D\) if necessary verifies every hypothesis
of Lemma~\ref{lem:linecover}.  Its conclusions give the first two
properties.

To prove the third, fix \(p\notin\ell\).  For one candidate in
group \(i\), there are at most \(s-1\) possible points
\(y\in\ell\) whose join to \(p\) has a direction in \(S_i\):
each such direction contributes at most one intersection with
\(\ell\).  A fixed point belongs to a uniform translate of
\(W_i\) with probability \(s^2/q^2\).  The probability for this
candidate is thus at most
\[
 \frac{(s-1)s^2}{q^2}.
\]
If six distinct groups meet the condition, choose one witnessing
candidate from each.  There are at most \((KN_q)^6=O_{s,K}(q^6)\)
choices, and their shifts are independent.  For the fixed pair
\((p,\ell)\), the probability of six such groups is therefore
\(O_{s,K}(q^{-6})\).  There are at most
\(q^2q(q+1)=O(q^4)\) pairs \((p,\ell)\).  The union bound gives
a total failure probability \(O_{s,K}(q^{-2})\), proving the claim.
\end{proof}

\subsection{Selection and transfer of the deletion estimates}

The estimates just proved concern independent candidates.  We can now
select compatible copies, using only inclusions to transfer the
estimates to the selected deletion set.
By Theorem~\ref{thm:selection}, with probability tending to one the
same pool admits one candidate per group with no conflicts.
Intersect this event with those of Lemma~\ref{con-pool-properties}.
No independence between these events is needed.  For every
sufficiently large odd integer \(n\), their intersection has positive
probability, so fix a pool in the intersection and a conflict-free
selection.  Denote the selected sets and shifts by
\[
 P_i=b_i+M_i(\F_s^2),\qquad Y_i=b_i+M_iY.
\]
Conflict avoidance gives the useful implication
\begin{equation}
\label{con-separation}
 i\ne j,\quad \operatorname{dir}(\ell)\in D_i,\quad
 \ell\cap P_i\ne\varnothing
 \quad\Longrightarrow\quad \ell\cap P_j=\varnothing.
\end{equation}
In particular the \(P_i\) are pairwise disjoint: a common point
would give a zero difference, which belongs to every conflict
direction, and each \(D_i\) is nonempty.

Let \(B\) be the union, as a set of lines, of the deletion lines
of the selected candidates, and put
\begin{equation}
\label{con-selected-survivors}
 X=\F_q^2\setminus\bigcup_{\ell\in B}\ell.
\end{equation}
Then \(X\supseteq X_{\rm pool}\).  Moreover, if \(\ell\notin B\),
every selected deletion line belongs to the pool list after all
exact occurrences of \(\ell\) have been omitted.  Hence \(X\)
contains the survivor set of that modified pool process. This
inclusion is deterministic: it holds even if \(\ell\) occurs in an
unselected candidate and even though the selected shifts depend on the
whole pool. In particular,
\begin{equation}
\label{con-selected-line-bound}
 \ell\notin B
 \quad\Longrightarrow\quad |X\cap\ell|\ge\frac{\delta q}{2}.
\end{equation}
The third property of Lemma~\ref{con-pool-properties} also remains
available for the selection, since it was established for all
candidates in the pool.

This completes the placement: the copies are separated, covers of the
ordinary survivors by at most \(q\) lines contain an almost pencil,
every undeleted line contains at least \(\delta q/2\) survivors,
and the bound of five groups remains valid.  The order of choices is
\(s,Y\), then the selection constant and \(K\), then the deletion
constants \(D,\delta,\rho\), and finally \(q=s^n\).
Only after fixing these constants do we require \(q=s^n\) to be
sufficiently large.  All further numerical inequalities below involve
these fixed constants
and can be imposed by enlarging this final threshold.  Thus the
construction is available for every sufficiently large odd integer
\(n\), with the same \(s,Y\); its pools and selected shifts may
depend on \(n\). The constant used for the ambient deletion lemma
may depend on \(K\), but it changes no hypothesis of the already
applied selection theorem. A small value of \(\delta\) only increases
the final lower threshold on \(q\).

\section{The hypergraph and its covering number}
\label{sec:hypergraph}

The placement supplies the ordinary points \(X\) and the exceptional
copies \(Y_i\).  We now turn them into edges.  Ordinary edges retain
all their affine-line labels.  In an exceptional copy, the directions
\(S_i\) instead use line labels read from \(f(Y)\); all remaining
directions retain their ordinary labels.  We will first check that
every pair of edges still meets, then use the rigidity of \(f(Y)\)
to rule out covers of size \(q\).

\subsection{Vertices, edges, and intersection}

Set \(r=q+1\).  There is one vertex part \(V_d\) for each
\(d\in\PP^1(\F_q)\).  For every ambient affine line \(\ell\)
of direction \(d\), place an ordinary vertex \(o_\ell\) in
\(V_d\).  If \(d\in S_i\), let \(t_i(d)\in\F_s^*\) be
the unique subfield direction whose image under \(M_i\) is
\(d\).  For each of the \(s\) subfield affine lines \(L\) of
direction \(t_i(d)\), place a special vertex \(v_{i,d,L}\) in
\(V_d\).  These subfield lines are read in the new coordinates
used for \(f(Y)\).  All vertices are formal labels, tagged to be
distinct between parts, between ordinary and special types, and
between different groups.  In particular each \(V_d\) is finite,
and the parts are pairwise disjoint.

Write \(\ell_d(z)\) for the ambient line of direction \(d\)
through \(z\), and \(L_t(w)\) for the subfield line of direction
\(t\) through \(w\in\F_s^2\).  Each \(x\in X\) defines the
ordinary edge
\[
 e_x=\{o_{\ell_d(x)}:d\in\PP^1(\F_q)\}.
\]
For each \(i\in\mathcal I\) and \(w\in Y\), put
\(y_{i,w}=b_i+M_iw\) and define an exceptional edge by selecting,
in part \(V_d\), the vertex
\begin{equation}
\label{con-exceptional-vertex}
 \begin{cases}
 v_{i,d,L_{t_i(d)}(f(w))},& d\in S_i,\\
 o_{\ell_d(y_{i,w})},&d\notin S_i.
 \end{cases}
\end{equation}
Denote this edge by \(e_{i,w}\), and let
\[
 H=\{e_x:x\in X\}\cup
   \{e_{i,w}:i\in\mathcal I,\ w\in Y\}.
\]
Thus an exceptional edge changes labels in the assigned directions,
without introducing any new part.  Every edge contains exactly one
vertex from every part, so \(H\)
is \((q+1)\)-partite and \((q+1)\)-uniform.  It is nonempty,
because \(Y\ne\varnothing\).

Figure~\ref{fig:two-coordinates} shows how to read an exceptional
edge from its parameter \(w\).  The deletion lines were defined from
its ordinary base point \(y_{i,w}\), while its special labels use
\(f(w)\).  The two axis parallel classes are preserved by \(f\),
which joins these descriptions in the intersection proof.

\begin{figure}[htbp]
\centering
\begin{tikzpicture}[
  >=Latex,
  point/.style={draw,rounded corners=2pt,align=center,inner sep=7pt},
  labelbox/.style={draw,align=center,text width=5.1cm,inner sep=7pt},
  every node/.style={font=\small}]
\node[point] (w) at (0,0) {$w=(w_1,w_2)\in Y$};
\node[point] (old) at (-3.5,-1.55)
  {Ambient coordinates\\[3pt]$y_{i,w}=b_i+M_iw\in Y_i$};
\node[point] (new) at (3.5,-1.55)
  {Small-plane coordinates\\[3pt]$f(w)=(w_1^3,w_2)\in f(Y)$};
\node[labelbox] (ordinary) at (-3.5,-3.65)
  {Ordinary vertex in $V_d$\\[3pt]$o_{\ell_d(y_{i,w})}$};
\node[labelbox] (special) at (3.5,-3.65)
  {Special vertex in $V_d$\\[3pt]$v_{i,d,L_{t_i(d)}(f(w))}$};
\draw[->] (w) -- node[above left,inner sep=2pt,xshift=-4pt,yshift=-2pt] {embed} (old);
\draw[->] (w) -- node[above right,inner sep=2pt,xshift=4pt,yshift=-2pt] {apply $f$} (new);
\draw[->] (old) -- node[left] {$d\notin S_i$} (ordinary);
\draw[->] (new) -- node[right] {$d\in S_i$} (special);
\end{tikzpicture}

\caption{The two label rules for the vertex of \(e_{i,w}\) in part
\(V_d\).  The left branch uses the embedded point \(b_i+M_iw\);
the right branch uses \(f(w)=(w_1^3,w_2)\) in the small plane.
For each \(d\in A_i\), two parameters give the same ordinary label
exactly when their images under \(f\) lie on the same corresponding
horizontal or vertical line.
The diagram records maps and label choices, rather than geometric
positions in the finite plane.}
\label{fig:two-coordinates}
\end{figure}

\begin{proposition}
\label{con-intersection}
The hypergraph \(H\) is intersecting.  Furthermore, all the edges
in its displayed parametrization are distinct.
\end{proposition}

\begin{proof}
There are four types of pairs, according to whether their edges are
ordinary or exceptional and, in the latter case, whether their group
indices agree.

Two ordinary edges with distinct base points share the ordinary
label of their join.  Next, an exceptional base point \(y\in Y_i\)
lies on a deletion line in every direction of the nonempty set
\(S_i\), so \(y\notin X\).  Its join with \(x\in X\) cannot
have direction in \(S_i\): otherwise that join would be a selected
deletion line containing \(x\).  Its ordinary label therefore
belongs to both the exceptional edge and \(e_x\).

For exceptional edges from distinct groups \(i,j\), the base
points are distinct by \eqref{con-separation}.  That same separation
condition ensures that their join has direction outside both
\(D_i\) and \(D_j\), and hence outside \(S_i\cup S_j\).
The two edges share its ordinary label.

Finally, consider \(e_{i,w}\) and \(e_{i,w'}\) with \(w\ne w'\).
The transformed points \(f(w),f(w')\) are distinct.  If their
join has finite nonzero slope \(t\), the exceptional edges share
the special label of this subfield line in the part of direction
\(M_i(t)\).  If their join is horizontal or vertical, the relevant
original coordinates are equal as well, because \(f\) preserves
these two parallel classes bijectively.  The original base points
then share an ordinary axis-direction label, in a direction of
\(A_i\).  This completes the intersection proof.

We also check that the parameters count distinct edges, as required
when we count exceptional edges below.  Two ordinary labels in
nonparallel directions determine the ordinary base point, so ordinary
edges are distinct.
An exceptional edge has special labels and cannot equal an ordinary
edge.  Within one group, the special labels in any two distinct
directions of \(S_i\) determine \(f(w)\), since the corresponding
subfield lines are nonparallel.  As \(|S_i|=s-1\ge2\), equal
exceptional edges within a group would imply \(w=w'\).  Finally,
any special label tagged by group \(i\) is absent from every edge
of a different group.  Exceptional edges in distinct groups are
therefore distinct too.
\end{proof}

\Needspace{14\baselineskip}
\subsection{Covering vertices in the small-plane coordinates}

We identify the positions of the exceptional edges of group \(i\)
with \(f(Y)\), via \(e_{i,w}\leftrightarrow f(w)\).  The action
of a covering vertex means the subset of these positions whose edges
contain that vertex.  The next lemma translates every such action
into one of the two resources in Lemma~\ref{lem:gadget}: a subfield
line or an individual point.

\begin{lemma}
\label{con-label-action}
Fix a group \(i\).
\begin{enumerate}
\item A special vertex of group \(i\) covers the positions on one
subfield line.  A special vertex of another group covers none.
\item An ordinary vertex of direction in \(S_i\) covers none.
An ordinary vertex of direction in \(A_i\), if it covers any
positions, covers them on a single horizontal or vertical line in
the \(f\) coordinates.
\item An ordinary vertex of direction outside \(S_i\cup A_i\)
covers at most one position.
\end{enumerate}
In particular, the positions covered by any one vertex are contained
in a subfield affine line.
\end{lemma}

\begin{proof}
The assertions for special vertices and altered ordinary directions
follow directly from \eqref{con-exceptional-vertex}.  For an axis
direction, undoing \(b_i,M_i\) fixes one original coordinate, so
applying \(f\) gives the corresponding horizontal or vertical
line.  A line in a direction outside \(S_i\cup A_i\) cannot
contain two points of \(P_i\), since the join directions of that
embedded subplane are exactly \(S_i\cup A_i\).  It therefore
covers at most one exceptional position.  A singleton can be
included in an arbitrary subfield line through it, and an empty
set requires no covering resource.
\end{proof}

\subsection{The final cover argument}

\begin{proof}[Completion of the proof of Theorem~\ref{thm:main}]
We prove that \(H\) has no vertex cover of size at most \(q\).
Suppose to the contrary that \(C\) is such a cover.  Its ordinary
vertices, regarded as affine lines, must cover \(X\): special
vertices belong to no ordinary edge.  These lines also cover
\(X_{\rm pool}\).  Lemma~\ref{con-pool-properties} therefore
gives a projective point \(p\) such that at least
\((1-\rho)q\) ordinary vertices of \(C\) are lines through
\(p\).  The line-survival bound will force every undeleted member
of this pencil into \(C\).  We then count the vertices left for
exceptional edges, separating centers at infinity, centers inside one
of the sets \(P_i\), and finite centers outside them all.

\paragraph{Every undeleted line in this pencil is forced.}
We claim that
\begin{equation}
\label{con-forced-pencil}
 p\in\ell,\quad\ell\notin B
 \quad\Longrightarrow\quad o_\ell\in C,
\end{equation}
where incidence with an infinite \(p\) means that \(\ell\) has
the corresponding direction.  If \(o_\ell\) were missing, all
other lines of \(C\) through \(p\) would cover at most one
affine point of \(\ell\).  There are at most \(\rho q\)
vertices outside the large pencil subcollection.  Each of their
ordinary lines covers at most one point of \(\ell\), and their
special vertices cover no ordinary edge.  Consequently
\[
 |X\cap\ell|\le1+\rho q.
\]
This contradicts \eqref{con-selected-line-bound}, since
\(\delta q/2>1+\rho q\) for sufficiently large \(q\) and
\(\rho=\delta/10\).  At infinity the pencil lines are parallel,
so the possible contribution of one central point is absent; the
same upper bound still applies.

Let \(F_p\) be the set of ordinary vertices forced by
\eqref{con-forced-pencil}, and call the vertices of \(C\setminus F_p\)
the additional vertices.  The deleted pencil lines determine how
many additional vertices the assumed budget \(|C|\le q\) permits.

\paragraph{An infinite center.}
Suppose \(p\) is the point at infinity in direction \(d\).  Put
\[
 I_d=\{i\in\mathcal I:d\in S_i\},
 \qquad z=|I_d|\in\{1,2\}.
\]
Only these groups contribute deletion lines in direction \(d\),
and each contributes \(s-1\), by
Lemma~\ref{con-deletion-count}.  At most \(z(s-1)\) of the
\(q\) lines in this parallel class belong to \(B\).  Thus
\[
 |C\setminus F_p|\le z(s-1).
\]
None of the forced vertices covers any exceptional edge of a group
in \(I_d\), because those groups use special vertices in
direction \(d\).

If \(z=1\), all positions of this group's \(f(Y)\) must therefore
be covered by at most \(s-1\) additional vertices.  By
Lemma~\ref{con-label-action}, their effects can be included in at
most \(s-1\) subfield lines.  This contradicts
Lemma~\ref{lem:gadget}.

If \(z=2\), the two groups are \(1,*\), by
\eqref{con-double-directions}.  An additional vertex can now meet
both groups, so we use the weighted line-and-point budget instead of
counting one line per vertex.  For each of these groups let \(a_i\)
count the additional vertices that cover at least one of its
positions and are either special vertices of this group or ordinary
vertices in its axis directions \(A_i\).  Let \(b_i'\) count the
other additional vertices covering a position of the group.  By
Lemma~\ref{con-label-action}, the latter vertices each cover only
a singleton.  The former supply subfield lines.  The weighted
conclusion of Lemma~\ref{lem:gadget} gives
\begin{equation}
\label{con-two-group-lower}
 a_i+\frac{b_i'}2\ge s
 \qquad (i=1,*).
\end{equation}

Each additional vertex contributes at most one to the sum of the
two left sides.  A special vertex is confined to its tagged group.
An ordinary axis-direction vertex counted in \(a_i\) is a line
meeting \(P_i\) with direction in \(A_i\subseteq D_i\); by
\eqref{con-separation} it cannot meet the other group's \(P_j\).
It therefore contributes nothing to that other group, even as a
singleton.  Any vertex of neither exclusive type contributes at
most \(1/2\) to each group.  This accounts for every type of
additional vertex, including those covering neither group.
Summing \eqref{con-two-group-lower}, we obtain
\[
 2s\le a_1+\frac{b_1'}2+a_*+\frac{b_*'}2
 \le |C\setminus F_p|\le2(s-1),
\]
a contradiction.

\paragraph{A finite center inside an embedded configuration.}
Now let \(p\) be finite, and write
\[
 h_p=|\{\ell\in B:p\in\ell\}|.
\]
There are \(q+1\) affine lines through \(p\), so
\begin{equation}
\label{con-finite-budget}
 |F_p|=q+1-h_p,
 \qquad |C\setminus F_p|\le h_p-1.
\end{equation}
The loss of one in the second inequality comes from the \(q+1\)
members of a finite pencil and the assumed cover size \(q\).
If \(h_p=0\), the forced vertices already exceed the permitted
size of \(C\).  We may therefore assume \(h_p\ge1\) in both
finite-center cases.

Suppose \(p\in P_i\), for the necessarily unique such group.
Every deletion line through \(p\) comes from this group: a line
from another group would contradict \eqref{con-separation}, since
its direction belongs to that other group's conflict set.  There
is at most one such line in each direction of \(S_i\), and hence
\begin{equation}
\label{con-internal-pencil-size}
 h_p\le s-1.
\end{equation}
Let \(w_p=M_i^{-1}(p-b_i)\in\F_s^2\), whether or not
\(w_p\in Y\).  We account for the exceptional positions of this
group covered by \(F_p\): they will all lie on two nonparallel
lines in the new coordinates.  For a position with base point
\(y\ne p\), its join with \(p\) has a direction in
\(S_i\cup A_i\).  A direction in \(S_i\) cannot cover this
edge by an ordinary vertex.  Any actual forced coverage must
therefore lie along one of the two axis-direction lines through
\(p\).  A position based at \(p\) itself, if present, lies on
both these axes as well.

Since \(f\) preserves the two coordinate parallel classes, all
these covered positions are contained, in the new coordinates, in
the horizontal and vertical lines through \(f(w_p)\).  The
remaining positions are covered by the additional vertices; each
can be included in one subfield line by
Lemma~\ref{con-label-action}.  Equations
\eqref{con-finite-budget} and \eqref{con-internal-pencil-size}
would thus give a cover of \(f(Y)\) by at most
\[
 2+(h_p-1)\le s
\]
lines containing the two nonparallel coordinate lines. These
are legitimate covering resources whether or not their ordinary labels
belong to \(F_p\): their union contains every position covered by
\(F_p\). Removing any duplicate replacement lines preserves these
two distinct lines and can only decrease the count.  This contradicts
Lemma~\ref{lem:gadget}.

\paragraph{A finite center outside all embedded configurations.}
It remains to consider \(p\notin\bigcup_{i\in\mathcal I}P_i\).
Here we use the five points on each deletion line, rather than the
small-plane cover obstruction.  We will find \(5h_p\) exceptional
edges missed by \(F_p\), while each additional vertex covers at
most five of them.

Each group contributes at most one deletion line through \(p\).
Indeed, two distinct such lines cannot be parallel.  After undoing
\(b_i,M_i\), their subfield affine equations have distinct
subfield directions.  Solving these two equations uses a
nonsingular linear system over \(\F_s\), so their intersection
lies in \(\F_s^2\).  This would place \(p\) in \(P_i\), a
contradiction.  Also, deletion lines from different groups cannot
coincide, by \eqref{con-separation}.

The \(h_p\) deletion lines through \(p\) therefore come from
\(h_p\) distinct groups.  From each such group \(i\), choose five
distinct points of \(Y_i\) on its contributing line, possible by
Lemma~\ref{con-deletion-count}.  These give \(5h_p\) distinct
exceptional edges, by Proposition~\ref{con-intersection}.  None
is covered by a forced ordinary vertex: its unique join line to
\(p\) is precisely its contributing deletion line, whose
direction is altered for that group.

Any one additional vertex covers at most five of these selected
edges.  For a special vertex this follows because it belongs to
only one group.  An ordinary line through \(p\) covers none,
by the same unique-join and altered-direction argument.  Finally,
let an ordinary vertex \(o_\ell\), with \(\ell\) not through
\(p\), cover some selected edges.  The line \(\ell\) meets each
contributing line in at most one point, so \(o_\ell\) covers at most
one selected edge from each contributing group.  Every such group
has a selected point \(y\in Y_i\cap\ell\) whose join to \(p\)
has direction in \(S_i\).  The third part of
Lemma~\ref{con-pool-properties}, which applies to all candidates
and hence to the selection, bounds the number of these groups by
five.  This proves the asserted bound for every additional vertex.

By \eqref{con-finite-budget}, at most \(h_p-1\) additional
vertices are available.  They can cover at most \(5(h_p-1)\) of
the \(5h_p\) selected edges, a final contradiction.

We have ruled out every possible projective center \(p\), and
hence every vertex cover of size at most \(q\).  Therefore
\(\tau(H)\ge q+1\).  Conversely, any edge of the nonempty
intersecting hypergraph \(H\) is itself a vertex cover and has
size \(q+1\).  It follows that
\[
 \tau(H)=q+1.
\]
The parameter choices in Section~\ref{sec:placement} give this
hypergraph for every sufficiently large odd integer \(n\), with
\(q=s^n\), after fixing any sufficiently large prime
\(s\equiv2\pmod3\) and a successful \(Y\). Choose one such \(Y\)
for each eligible \(s\), and let \(n_0(s)\ge3\) exceed all the
resulting degree thresholds. This proves Theorem~\ref{thm:main}.
Since there are arbitrarily large odd integers, it also gives infinitely
many ranks \(r=q+1\).
\end{proof}
